\documentclass[11pt,a4paper]{article}

% Packages
\usepackage[utf8]{inputenc}
\usepackage[T1]{fontenc}
\usepackage{geometry}
\usepackage{listings}
\usepackage{xcolor}
\usepackage{hyperref}
\usepackage{fancyhdr}
\usepackage{tcolorbox}
\usepackage{enumitem}
\usepackage{graphicx}
\usepackage{booktabs}

% Page geometry
\geometry{margin=1in}

% Colors
\definecolor{codegreen}{rgb}{0,0.6,0}
\definecolor{codegray}{rgb}{0.5,0.5,0.5}
\definecolor{backcolour}{rgb}{0.95,0.95,0.92}
\definecolor{linkblue}{rgb}{0.1,0.3,0.6}

% Code listings style
\lstdefinestyle{bashstyle}{
    backgroundcolor=\color{backcolour},
    basicstyle=\ttfamily\small,
    breaklines=true,
    frame=single,
    framerule=0.5pt,
    numbers=none,
    commentstyle=\color{codegreen},
}
\lstset{style=bashstyle}

% Header/Footer
\pagestyle{fancy}
\fancyhf{}
\rhead{PIQC User Guide}
\lhead{v1.0.0}
\rfoot{Page \thepage}

% Hyperref setup
\hypersetup{
    colorlinks=true,
    linkcolor=linkblue,
    urlcolor=linkblue,
    pdftitle={PIQC User Guide},
    pdfauthor={ParallelIQ},
}

\begin{document}

% Title Page
\begin{center}
    {\Huge\textbf{PIQC User Guide}}\\[0.5em]
    {\Large Production Inference Quality Control}\\[2em]
    {\large Version 1.0.0}\\[1em]
    \today
\end{center}

\vspace{2em}

\begin{tcolorbox}[colback=blue!5,colframe=blue!50!black,title=\textbf{What is PIQC?}]
\textbf{PIQC (Production Inference Quality Control)} is a Kubernetes-native tool that automatically discovers AI/ML inference deployments, measures their efficiency, and surfaces the dollar cost of idle and unallocated GPUs.

\textbf{Key Features:}
\begin{itemize}
    \item Automatic vLLM deployment discovery across all namespaces
    \item GPU metrics collection (utilization, memory, temperature) via \texttt{nvidia-smi}
    \item Runtime metrics from vLLM API (latency, throughput, KV cache)
    \item \textbf{MFU (Model FLOPS Utilization)} per deployment
    \item \textbf{Cost-per-1K-tokens} — translate GPU spend into business metrics
    \item \textbf{Revenue leak reporting} — idle GPU waste and unallocated node detection
    \item Default table output: run \texttt{piqc scan} and see dollar figures immediately
    \item Works remotely via kubeconfig or in-cluster as a Kubernetes Job
\end{itemize}
\end{tcolorbox}

\newpage
\tableofcontents
\newpage

%==============================================================================
\section{Installation}
%==============================================================================

\subsection{Prerequisites}

\begin{itemize}
    \item Python 3.11 or higher
    \item \texttt{kubectl} configured with cluster access
    \item \texttt{poetry} (Python dependency manager)
    \item Access to a Kubernetes cluster with vLLM deployments
\end{itemize}

%------------------------------------------------------------------------------
\subsection{Download from GitHub}

\begin{lstlisting}[language=bash]
# Clone the repository
git clone https://github.com/your-org/ModelSpec.git
cd ModelSpec

# Or download ZIP and extract
# wget https://github.com/your-org/ModelSpec/archive/main.zip
# unzip main.zip
# cd ModelSpec-main
\end{lstlisting}

%------------------------------------------------------------------------------
\subsection{Install Dependencies}

PIQC uses Poetry for dependency management:

\begin{lstlisting}[language=bash]
# Install Poetry if not already installed
curl -sSL https://install.python-poetry.org | python3 -

# Install PIQC dependencies
poetry install

# Verify installation
poetry run piqc version
\end{lstlisting}

\textbf{Expected output:}
\begin{lstlisting}
ModelSpec Introspector v1.0.0
\end{lstlisting}

%==============================================================================
\section{Quick Start: Testing vLLM Deployments}
%==============================================================================

\subsection{Verify Cluster Access}

\begin{lstlisting}[language=bash]
# Test Kubernetes connection
poetry run piqc test-connection
\end{lstlisting}

\textbf{Expected output:}
\begin{lstlisting}
[INFO] Connecting to cluster...
       Context: my-k8s-cluster
       Cluster: my-k8s-cluster
[INFO] Connection successful!
\end{lstlisting}

%------------------------------------------------------------------------------
\subsection{Scan for vLLM Deployments}

\subsubsection{Basic Scan — Cost Report (Default)}

Running \texttt{piqc scan} with no flags connects to your current kubectl context, scans all namespaces, and immediately prints a cost report:

\begin{lstlisting}[language=bash]
poetry run piqc scan
\end{lstlisting}

The output includes GPU utilization, MFU, \$/1K tokens, idle waste per deployment, and a Cost Summary panel showing total leaked revenue per day and annualized.

\subsubsection{Scan Specific Namespace}

\begin{lstlisting}[language=bash]
poetry run piqc scan -n production
\end{lstlisting}

\subsubsection{Override GPU Cost Rate}

If your actual cloud or lease rate differs from the built-in estimates:

\begin{lstlisting}[language=bash]
# Set GPU cost to $3.50/GPU/hr
poetry run piqc scan --gpu-cost 3.50

# Set separate rate for unallocated (dark) nodes
poetry run piqc scan --gpu-cost 3.50 --node-cost 3.50
\end{lstlisting}

\subsubsection{Full YAML Output}

\begin{lstlisting}[language=bash]
# Generate detailed ModelSpec YAML files
poetry run piqc scan --format yaml -o ./output
\end{lstlisting}

Files are saved to \texttt{./output/} by default.

%------------------------------------------------------------------------------
\subsection{Runtime Metrics}

Runtime metrics (token throughput, latency, KV cache) are collected automatically on every scan via the vLLM API. They are required for MFU and \$/1K tokens calculations.

To disable runtime collection (faster scan, no MFU or cost-per-token):

\begin{lstlisting}[language=bash]
poetry run piqc scan --no-collect-runtime
\end{lstlisting}

\textbf{What runtime collection provides:}
\begin{itemize}
    \item Request latency percentiles (TTFT, TPOT, E2E)
    \item Token throughput (prompt \& generation tokens/sec)
    \item KV cache and prefix cache hit rates
    \item MFU (Model FLOPS Utilization)
    \item Cost per 1K generated tokens
\end{itemize}

%==============================================================================
\section{Testing on Google Cloud Platform (GCP)}
%==============================================================================

\subsection{Option 1: Using Cloud Shell (Recommended)}

Cloud Shell provides a pre-configured environment with \texttt{kubectl} and \texttt{python}.

\subsubsection{Step 1: Open Cloud Shell}

\begin{enumerate}
    \item Go to \url{https://console.cloud.google.com}
    \item Click the \textbf{Activate Cloud Shell} icon (top-right)
    \item Wait for the terminal to load
\end{enumerate}

\subsubsection{Step 2: Set Up Your Project}

\begin{lstlisting}[language=bash]
# Check current project
gcloud config get-value project

# Set project if needed
gcloud config set project YOUR_PROJECT_ID

# Enable required APIs
gcloud services enable container.googleapis.com
gcloud services enable compute.googleapis.com
\end{lstlisting}

\subsubsection{Step 3: Create GKE Cluster with GPU}

\begin{lstlisting}[language=bash,basicstyle=\ttfamily\footnotesize]
gcloud container clusters create piqc-test \
  --zone=us-central1-a \
  --machine-type=n1-standard-4 \
  --accelerator=type=nvidia-tesla-t4,count=1 \
  --num-nodes=1 \
  --spot \
  --disk-size=50GB
\end{lstlisting}

\textbf{Cost:} \textasciitilde\$0.50/hour with Spot instances

\subsubsection{Step 4: Install GPU Drivers}

\begin{lstlisting}[language=bash,basicstyle=\ttfamily\footnotesize]
kubectl apply -f https://raw.githubusercontent.com/GoogleCloudPlatform/container-engine-accelerators/master/nvidia-driver-installer/cos/daemonset-preloaded.yaml

# Wait for drivers to install (~2 minutes)
kubectl get pods -n kube-system -w | grep nvidia
\end{lstlisting}

\subsubsection{Step 5: Deploy Test vLLM Server}

\begin{lstlisting}[language=bash]
# Create namespace
kubectl create namespace inference

# Deploy TinyLlama on vLLM
cat << 'EOF' | kubectl apply -f -
apiVersion: apps/v1
kind: Deployment
metadata:
  name: vllm-tinyllama
  namespace: inference
  labels:
    app: vllm-tinyllama
    app.kubernetes.io/name: vllm
    framework: vllm
spec:
  replicas: 1
  selector:
    matchLabels:
      app: vllm-tinyllama
  template:
    metadata:
      labels:
        app: vllm-tinyllama
        app.kubernetes.io/name: vllm
        framework: vllm
    spec:
      containers:
      - name: vllm
        image: vllm/vllm-openai:v0.6.3.post1
        args:
          - "--model"
          - "TinyLlama/TinyLlama-1.1B-Chat-v1.0"
          - "--served-model-name"
          - "tinyllama"
          - "--dtype"
          - "float16"
          - "--max-model-len"
          - "1024"
          - "--gpu-memory-utilization"
          - "0.5"
        ports:
          - containerPort: 8000
        resources:
          requests:
            nvidia.com/gpu: 1
            cpu: "2"
            memory: "8Gi"
          limits:
            nvidia.com/gpu: 1
            cpu: "4"
            memory: "12Gi"
---
apiVersion: v1
kind: Service
metadata:
  name: vllm-tinyllama
  namespace: inference
spec:
  selector:
    app: vllm-tinyllama
  ports:
    - port: 8000
      targetPort: 8000
  type: ClusterIP
EOF

# Wait for pod to be ready (~3-5 minutes)
kubectl get pods -n inference -w
\end{lstlisting}

\subsubsection{Step 6: Upload PIQC to Cloud Shell}

\textbf{Method A: Clone from GitHub}
\begin{lstlisting}[language=bash]
cd ~
git clone https://github.com/your-org/ModelSpec.git
cd ModelSpec
poetry install
\end{lstlisting}

\textbf{Method B: Upload ZIP File}
\begin{enumerate}
    \item Create ZIP on your local machine: \texttt{ModelSpec.zip}
    \item In Cloud Shell, click \textbf{⋮} menu → \textbf{Upload}
    \item Select \texttt{ModelSpec.zip}
    \item Extract and install:
\end{enumerate}

\begin{lstlisting}[language=bash]
cd ~
unzip ModelSpec.zip
cd ModelSpec
poetry install
\end{lstlisting}

\subsubsection{Step 7: Apply RBAC Permissions}

\begin{lstlisting}[language=bash]
cd ~/ModelSpec
kubectl apply -f rbac/
\end{lstlisting}

This grants PIQC permission to exec into pods and list cluster resources.

\subsubsection{Step 8: Run PIQC Tests}

\begin{lstlisting}[language=bash]
# Basic scan
poetry run piqc scan --namespace inference --format table

# Full scan with runtime metrics
poetry run piqc scan --namespace inference \
    --collect-runtime \
    --format yaml \
    -o ~/piqc-outputs

# View generated file
cat ~/piqc-outputs/inference_vllm-tinyllama.yaml
\end{lstlisting}

\subsubsection{Step 9: Download Results}

\begin{lstlisting}[language=bash]
# Create ZIP of outputs
cd ~
zip -r piqc-outputs.zip piqc-outputs/

# Download via Cloud Shell menu
# Click ⋮ → Download → enter "piqc-outputs.zip"
\end{lstlisting}

\subsubsection{Step 10: Cleanup (Important!)}

\begin{lstlisting}[language=bash]
# Delete vLLM deployment
kubectl delete deployment vllm-tinyllama -n inference

# Delete cluster to stop charges
gcloud container clusters delete piqc-test \
  --zone=us-central1-a \
  --quiet
\end{lstlisting}

%==============================================================================
\section{Using PIQC from Your Local Machine}
%==============================================================================

PIQC can also run from your laptop/workstation using an existing kubeconfig.

\subsection{Prerequisites}

\begin{itemize}
    \item Python 3.11+
    \item \texttt{kubectl} installed and configured
    \item Valid kubeconfig with cluster access
    \item Poetry installed
\end{itemize}

%------------------------------------------------------------------------------
\subsection{Windows Setup}

\subsubsection{Install Python}

Download Python 3.11+ from \url{https://www.python.org/downloads/}

\subsubsection{Install kubectl}

\begin{lstlisting}[language=bash]
# Using Chocolatey
choco install kubernetes-cli

# Or download from https://kubernetes.io/docs/tasks/tools/
\end{lstlisting}

\subsubsection{Install Poetry}

\begin{lstlisting}[language=bash]
# PowerShell
(Invoke-WebRequest -Uri https://install.python-poetry.org -UseBasicParsing).Content | python -
\end{lstlisting}

Add to PATH: \texttt{C:\textbackslash Users\textbackslash <YourName>\textbackslash AppData\textbackslash Roaming\textbackslash Python\textbackslash Scripts}

\subsubsection{Configure kubectl}

Option 1: Download kubeconfig from GKE
\begin{lstlisting}[language=bash]
gcloud container clusters get-credentials CLUSTER_NAME \
  --zone=ZONE \
  --project=PROJECT_ID
\end{lstlisting}

Option 2: Use provided kubeconfig file
\begin{lstlisting}[language=bash]
# Save kubeconfig to specific location
$env:KUBECONFIG = "C:\path\to\kubeconfig.yaml"
kubectl get nodes
\end{lstlisting}

\subsubsection{Install and Run PIQC}

\begin{lstlisting}[language=bash]
# Clone repository
cd C:\Users\YourName\Desktop
git clone https://github.com/your-org/ModelSpec.git
cd ModelSpec

# Install dependencies
poetry install

# Test connection
poetry run piqc test-connection

# Run scan
poetry run piqc scan --namespace inference \
  --collect-runtime \
  --format yaml \
  -o C:\Users\YourName\Desktop\piqc-outputs
\end{lstlisting}

%------------------------------------------------------------------------------
\subsection{Linux/macOS Setup}

\subsubsection{Install Prerequisites}

\begin{lstlisting}[language=bash]
# macOS (using Homebrew)
brew install python@3.11
brew install kubectl

# Ubuntu/Debian
sudo apt update
sudo apt install python3.11 python3-pip kubectl

# Install Poetry
curl -sSL https://install.python-poetry.org | python3 -
\end{lstlisting}

\subsubsection{Configure kubectl}

\begin{lstlisting}[language=bash]
# Get credentials from GKE
gcloud container clusters get-credentials CLUSTER_NAME \
  --zone=ZONE \
  --project=PROJECT_ID

# Or set custom kubeconfig
export KUBECONFIG=/path/to/kubeconfig.yaml
kubectl get nodes
\end{lstlisting}

\subsubsection{Install and Run PIQC}

\begin{lstlisting}[language=bash]
# Clone repository
cd ~/Desktop
git clone https://github.com/your-org/ModelSpec.git
cd ModelSpec

# Install dependencies
poetry install

# Run scan
poetry run piqc scan --namespace inference \
  --collect-runtime \
  --format yaml \
  -o ~/piqc-outputs
\end{lstlisting}

%==============================================================================
\section{Command Reference}
%==============================================================================

\subsection{Core Commands}

\begin{table}[h]
\small
\begin{tabular}{@{}p{4cm}p{9cm}@{}}
\toprule
\textbf{Command} & \textbf{Description} \\
\midrule
\texttt{piqc version} & Show version information \\
\texttt{piqc test-connection} & Verify cluster connectivity \\
\texttt{piqc scan} & Scan cluster for vLLM deployments \\
\bottomrule
\end{tabular}
\end{table}

\subsection{Scan Options}

\begin{table}[h]
\small
\begin{tabular}{@{}p{5.5cm}p{7.5cm}@{}}
\toprule
\textbf{Option} & \textbf{Description} \\
\midrule
\texttt{--namespace, -n} & Scan specific namespace (default: all) \\
\texttt{--format} & Output format: \textbf{table} (default), yaml, json \\
\texttt{--output, -o} & Output directory for yaml/json files \\
\texttt{--gpu-cost DOLLARS} & Override GPU cost in \$/GPU/hr \\
\texttt{--node-cost DOLLARS} & Override cost for unallocated node GPUs \\
\texttt{--collect-runtime / --no-collect-runtime} & Runtime metrics via vLLM API (default: on) \\
\texttt{--no-exec} & Skip GPU metrics collection via pod exec \\
\texttt{--verbose, -v} & Enable verbose logging \\
\texttt{--debug} & Enable debug trace logging \\
\texttt{--workers} & Number of parallel scan workers \\
\texttt{--combined} & Output single combined file \\
\texttt{--output-piqc} & Generate PIQC facts bundle \\
\bottomrule
\end{tabular}
\end{table}

\subsection{Example Commands}

\begin{lstlisting}[language=bash,basicstyle=\ttfamily\footnotesize]
# Quick table view of all deployments
poetry run piqc scan --format table

# Detailed YAML for specific namespace
poetry run piqc scan --namespace prod --format yaml

# Full scan with runtime metrics
poetry run piqc scan --namespace inference \
  --collect-runtime \
  --format yaml \
  -o ./results

# Fast scan without GPU metrics
poetry run piqc scan --no-exec --format json

# Debug mode for troubleshooting
poetry run piqc scan --namespace test --debug

# Generate ParallelIQ facts bundle
poetry run piqc scan --namespace inference \
  --output-piqc \
  -o ./piqc-facts
\end{lstlisting}

%==============================================================================
\section{Output Formats}
%==============================================================================

\subsection{Table Format}

Default output — run \texttt{piqc scan} to see this immediately:

\begin{lstlisting}[basicstyle=\ttfamily\tiny]
┏━━━━━━━━━━━━━━━━━━━━━━━━━━━━━━┳━━━━━━━━━┳━━━━━━━━━━━━━━━━━━━┳━━━━━━━━━━┳━━━━━━━━━━┳━━━━━━━━┳━━━━━━━━━━━━┳━━━━━━━━┳━━━━━━━━━━━━━━┓
┃ Deployment                   ┃ Engine  ┃ GPU               ┃ Replicas ┃ GPU Util ┃ MFU    ┃ $/1K tokens┃ $/hr   ┃ Idle $/day   ┃
┡━━━━━━━━━━━━━━━━━━━━━━━━━━━━━━╇━━━━━━━━━╇━━━━━━━━━━━━━━━━━━━╇━━━━━━━━━━╇━━━━━━━━━━╇━━━━━━━━╇━━━━━━━━━━━━╇━━━━━━━━╇━━━━━━━━━━━━━━┩
│ meta-llama/Llama-3-70B-Inst  │ vllm    │ 8xA100-SXM4-80GB  │        2 │       4% │   3.1% │    $0.0842 │ $56.00 │   $1,290.24  │
│ mistral-7b-instruct          │ vllm    │ 1xA100-SXM4-40GB  │        1 │      11% │   8.4% │    $0.0156 │  $2.50 │      $53.40  │
│ codellama-34b-staging        │ vllm    │ 4xA100-SXM4-80GB  │        1 │       0% │   0.0% │        N/A │ $14.00 │     $336.00  │
│ embedding-bge-large          │ vllm    │ 1xT4              │        3 │      82% │  51.2% │    $0.0003 │  $1.35 │       $5.83  │
└──────────────────────────────┴─────────┴───────────────────┴──────────┴──────────┴────────┴────────────┴────────┴──────────────┘

Cost Summary
  Total GPU spend rate     : $73.85/hr
  Leased & idle (util <60%): $1,685.47/day  (pods running, GPUs underused)
  Unallocated nodes        :   $336.00/day  (4 GPUs with no pods scheduled)
  Total estimated leak     : $2,021.47/day  ($737,836/yr)
  Avg MFU (active)         : 15.7%  (healthy range: 30-60%)
\end{lstlisting}

\textbf{Column definitions:}
\begin{itemize}
    \item \textbf{MFU} — Model FLOPS Utilization: observed compute vs.\ theoretical GPU peak. Green $\geq$30\%, yellow $\geq$10\%, red $<$10\%.
    \item \textbf{\$/1K tokens} — Cost per 1,000 generated tokens based on live throughput and GPU spend rate.
    \item \textbf{Idle \$/day} — Estimated daily waste for this deployment. Red = utilization below 60\% threshold.
\end{itemize}

\subsection{YAML Format}

Complete ModelSpec with all collected data:

\begin{lstlisting}[language=yaml,basicstyle=\ttfamily\scriptsize]
apiVersion: modelspec.paralleliq.ai/v1
kind: ModelSpec
metadata:
  name: vllm-tinyllama
  namespace: inference
model:
  name: TinyLlama/TinyLlama-1.1B-Chat-v1.0
  architecture: llama
  parameters: 1.1B
resources:
  gpus:
  - type: Tesla T4
    memoryTotal: 15GB
    temperature: 55
runtimeState:
  vllm:
    loadedModel: tinyllama
    healthStatus: healthy
    promptTokensPerSec: 1250.5
\end{lstlisting}

\subsection{JSON Format}

Machine-readable format for automation:

\begin{lstlisting}[language=json,basicstyle=\ttfamily\scriptsize]
{
  "apiVersion": "modelspec.paralleliq.ai/v1",
  "kind": "ModelSpec",
  "metadata": {
    "name": "vllm-tinyllama",
    "namespace": "inference"
  },
  "model": {
    "name": "TinyLlama/TinyLlama-1.1B-Chat-v1.0"
  }
}
\end{lstlisting}

\subsection{PIQC Facts Bundle}

Standardized facts format for ParallelIQ platform:

\begin{lstlisting}[language=json,basicstyle=\ttfamily\scriptsize]
{
  "schema": "piqc-scan.v0.1",
  "workloadObjects": [
    {
      "id": "inference/vllm-tinyllama",
      "facts": {
        "runtime.engineType": "vllm",
        "model.name": "TinyLlama-1.1B-Chat-v1.0",
        "hardware.gpuType": "Tesla T4"
      }
    }
  ]
}
\end{lstlisting}

%==============================================================================
\section{Troubleshooting}
%==============================================================================

\subsection{Common Issues}

\subsubsection{Connection Errors}

\textbf{Error:} \texttt{Unable to connect to cluster}

\textbf{Solution:}
\begin{lstlisting}[language=bash]
# Verify kubectl works
kubectl get nodes

# Check current context
kubectl config current-context

# Switch context if needed
kubectl config use-context my-cluster
\end{lstlisting}

\subsubsection{RBAC Permission Errors}

\textbf{Error:} \texttt{User cannot list pods in namespace}

\textbf{Solution:}
\begin{lstlisting}[language=bash]
# Apply RBAC permissions
kubectl apply -f rbac/

# Or ask admin to grant cluster-viewer role
\end{lstlisting}

\subsubsection{GPU Metrics Not Available}

\textbf{Error:} \texttt{nvidia-smi: command not found}

\textbf{Solution:}
\begin{itemize}
    \item Install GPU drivers in cluster
    \item Or use \texttt{--no-exec} flag to skip GPU metrics
\end{itemize}

\subsubsection{Runtime Metrics Not Collected}

\textbf{Error:} \texttt{Could not discover vLLM service endpoint}

\textbf{Solution:}
\begin{itemize}
    \item Verify vLLM pod is running: \texttt{kubectl get pods -n <namespace>}
    \item Check vLLM is serving on port 8000
    \item Ensure RBAC allows pod exec
\end{itemize}

\subsection{Debug Mode}

For detailed troubleshooting:

\begin{lstlisting}[language=bash]
poetry run piqc scan --namespace inference --debug
\end{lstlisting}

\subsection{Getting Help}

\begin{itemize}
    \item GitHub Issues: \url{https://github.com/your-org/ModelSpec/issues}
    \item Documentation: \url{https://your-org.github.io/ModelSpec}
    \item Email: support@paralleliq.ai
\end{itemize}

%==============================================================================
\section{Best Practices}
%==============================================================================

\begin{enumerate}
    \item \textbf{Run \texttt{piqc scan} first}: The default table output gives an immediate cost report with no flags needed.
    \item \textbf{Set your actual GPU rate}: Use \texttt{--gpu-cost} with your real lease or cloud rate so the dollar figures are accurate, not estimates.
    \item \textbf{Prioritize red MFU rows}: Deployments with MFU below 10\% are strong decommission or consolidation candidates.
    \item \textbf{Act on unallocated nodes first}: These are GPUs with no pods at all — pure spend with zero utilization. Schedule workloads or return the capacity.
    \item \textbf{Use \texttt{--no-collect-runtime} for fast audits}: Skips vLLM API calls for a quicker scan when MFU and \$/1K tokens are not needed.
    \item \textbf{Use namespaces to scope}: Use \texttt{-n production} to focus on live workloads and exclude staging noise.
    \item \textbf{Save YAML for records}: Use \texttt{--format yaml -o ./snapshots} to build a history of deployment state.
    \item \textbf{Automate}: Run PIQC as a Kubernetes CronJob in-cluster to generate daily cost reports without manual intervention.
\end{enumerate}

\end{document}
